%=====================================================================
% Datos · Registros cargados
%=====================================================================
\subsection*{Registros cargados}
\addcontentsline{toc}{subsection}{Registros cargados}

Los registros con que se probó la base de datos son los que carga \at{bd/insertar\_datos\_prueba.sql}, el script de la sección anterior: \dpNumFilas{} registros repartidos en las \dpNumTablasConDatos{} tablas del modelo, muchos más que los veinte que se piden como mínimo. No son filas de relleno. Forman un escenario del juego, fotografiado el \dpFecha{} a las 18:00 UTC, en el que cada registro representa una situación que describen los casos de uso: cuentas en cada estado, partidas de los tres tipos terminadas de las cinco formas, compras, cajas, reportes y sanciones. Cada cosa que los datos deben permitir probar tiene su subsección:

\begin{center}\small
\begin{longtable}{|L{0.22\textwidth}|L{0.2\textwidth}|L{0.48\textwidth}|}
\hline
\textbf{Qué se prueba} & \textbf{Dónde} & \textbf{Con qué registros} \\ \hline
\endfirsthead
\hline
\textbf{Qué se prueba} & \textbf{Dónde} & \textbf{Con qué registros} \\ \hline
\endhead
Casos normales & \textit{Casos normales} & Una partida clasificatoria de principio a fin y las compras en la tienda, con todo lo que dejan escrito en otras tablas. \\ \hline
Relaciones entre tablas & \textit{Relaciones entre tablas} & Consultas que recorren las llaves foráneas: de la partida a sus jugadores, de la cuenta a lo que lleva equipado, del reporte a la sanción y a la apelación, de la caja a su contenido, y de una sanción a la que la sustituye. \\ \hline
Diferentes estados o categorías & \textit{Estados y categorías} & \dnNumValoresConFilas{} de los \dnNumValores{} valores posibles de las \dnNumDominios{} columnas de estado o de categoría tienen al menos un registro; los que faltan se explican uno por uno. \\ \hline
Variaciones de la información & \textit{Variaciones} & Nulos que dependen del tipo de fila, cuentas eliminadas y baneadas, textos en inglés y con emojis, partidas sin reloj, de cuatro jugadores y perdidas por tiempo, jugadas deshechas, objetos retirados. \\ \hline
Que la base cumple sus reglas & \textit{Casos que la base rechaza} & \dnNumPruebasRechazo{} operaciones que el documento prohíbe, intentadas sobre estos registros; la base rechaza las \dnNumPruebasRechazadas{}. \\ \hline
Todos los registros & \textit{Registros por tabla} & El contenido de cada tabla, al final de la sección. \\ \hline
\end{longtable}
\end{center}

\textbf{De dónde salen las tablas.} Ninguna tabla de registros de esta sección se copió a mano: son el resultado de las consultas de \at{bd/consultas\_datos.sql} sobre la base cargada, que \at{bd/exportar\_datos.py} ejecuta y convierte en tablas del documento. Las mismas consultas pueden abrirse en SQL Server Management Studio y ejecutarse una a una. Para regenerar las tablas después de cambiar el escenario:

\begin{lstlisting}[style=sql]
python bd\generar_datos_prueba.py
sqlcmd -S localhost -E -b -I -f 65001 -i bd\insertar_datos_prueba.sql
python bd\exportar_datos.py
\end{lstlisting}

\textbf{Las capturas.} Se tomaron el 13 de septiembre de 2026 en la consola de Windows, ejecutando \at{sqlcmd} sobre SQL Server 2025 (17.0.4065.4). Cada una muestra el comando y su salida tal como aparecieron; la consola no tiene fuente para el emoji de un mensaje de chat y lo dibuja como un recuadro, aunque la base lo guarda completo. La captura 1 es la carga: los tres scripts, en orden, sobre la instancia, y las \dpNumComprobaciones{} comprobaciones del final de \at{insertar\_datos\_prueba.sql}, todas con cero incoherencias.

\captura{01-carga}{Captura 1. Creación de la base y de las tablas, carga de los registros y comprobaciones de coherencia.}

\textbf{Registros por tabla.} La tabla siguiente, que genera el script de datos, cuenta los registros de cada tabla por área; la captura 2 los cuenta en la propia base, consultando el catálogo del sistema, y da las mismas cifras.

% Archivo generado por bd/generar_datos_prueba.py. No editar a mano.
\begin{center}\small
\begin{longtable}{|L{0.34\textwidth}|L{0.12\textwidth}|}
\hline
\textbf{Tabla} & \textbf{Filas} \\ \hline
\endfirsthead
\hline
\textbf{Tabla} & \textbf{Filas} \\ \hline
\endhead
\multicolumn{2}{|l|}{\textbf{Catálogos precargados (\mbox{D-09})}} \\ \hline
\textit{Ranura} & 6 \\ \hline
\textit{ObjetoCosmetico} & 23 \\ \hline
\textit{Modo} & 3 \\ \hline
\textit{Division} & 4 \\ \hline
\textit{NivelIA} & 4 \\ \hline
\textit{Leccion} & 7 \\ \hline
\textit{TipoCaja} & 3 \\ \hline
\textit{TipoCajaObjeto} & 14 \\ \hline
\textit{PalabraProhibida} & 8 \\ \hline
\textit{MotivoReporte} & 5 \\ \hline
\multicolumn{2}{|l|}{\textbf{Cuenta y acceso}} \\ \hline
\textit{Usuario} & 11 \\ \hline
\textit{Sesion} & 15 \\ \hline
\textit{CodigoSegundoFactor} & 4 \\ \hline
\textit{TokenVerificacionCorreo} & 12 \\ \hline
\textit{TokenRecuperacion} & 3 \\ \hline
\textit{AceptacionTerminos} & 11 \\ \hline
\multicolumn{2}{|l|}{\textbf{Perfil y personalización}} \\ \hline
\textit{HistorialNickname} & 1 \\ \hline
\textit{Avatar} & 3 \\ \hline
\textit{EnlaceRed} & 2 \\ \hline
\textit{UsuarioObjeto} & 95 \\ \hline
\textit{Equipamiento} & 66 \\ \hline
\multicolumn{2}{|l|}{\textbf{Partida}} \\ \hline
\textit{Partida} & 12 \\ \hline
\textit{Participacion} & 25 \\ \hline
\textit{ParticipacionObjeto} & 150 \\ \hline
\textit{Jugada} & 134 \\ \hline
\textit{OfertaTablas} & 4 \\ \hline
\textit{EnlaceEspectador} & 1 \\ \hline
\multicolumn{2}{|l|}{\textbf{Emparejamiento, salas y chat}} \\ \hline
\textit{ColaEmparejamiento} & 1 \\ \hline
\textit{Sala} & 2 \\ \hline
\textit{SalaParticipante} & 2 \\ \hline
\textit{Invitacion} & 4 \\ \hline
\textit{Mensaje} & 11 \\ \hline
\multicolumn{2}{|l|}{\textbf{Clasificación y economía}} \\ \hline
\textit{EstadisticaModo} & 33 \\ \hline
\textit{HistorialDivision} & 2 \\ \hline
\textit{Caja} & 3 \\ \hline
\textit{CajaContenido} & 3 \\ \hline
\textit{MovimientoMoneda} & 30 \\ \hline
\multicolumn{2}{|l|}{\textbf{Relaciones entre jugadores y tutorial}} \\ \hline
\textit{Amistad} & 4 \\ \hline
\textit{Solicitud} & 1 \\ \hline
\textit{Silencio} & 1 \\ \hline
\textit{Bloqueo} & 1 \\ \hline
\textit{ProgresoTutorial} & 8 \\ \hline
\multicolumn{2}{|l|}{\textbf{Moderación y bitácoras}} \\ \hline
\textit{Reporte} & 5 \\ \hline
\textit{Sancion} & 5 \\ \hline
\textit{Apelacion} & 3 \\ \hline
\textit{BitacoraModeracion} & 19 \\ \hline
\textit{BitacoraAcceso} & 54 \\ \hline
\textbf{Total} & \textbf{818} \\ \hline
\end{longtable}
\end{center}


\captura{02-conteo}{Captura 2. Registros de cada tabla según SQL Server.}
